\chapter{线性代数：相似理论}

\section*{考研解题总纲}
\addcontentsline{toc}{section}{考研解题总纲}

\begin{itemize}
  \item 相似矩阵有相同特征值、迹、行列式、秩和特征多项式。
  \item 可对角化的判定：每个特征值的几何重数等于代数重数，或存在 \(n\) 个线性无关特征向量。
  \item 实对称矩阵必可正交对角化，不同特征值对应特征向量正交。
  \item Jordan 或相似标准形题要先求特征多项式，再求各特征值的特征子空间维数。
\end{itemize}



\section{第 31 页}
\begin{problemblock}
\begin{center}
\includegraphics[width=.94\textwidth]{figures/xiandai/page_031.png}
\end{center}
\end{problemblock}

\begin{solutionblock}
\textbf{本页考点定位：}第 31 页是本章第 1 页，主要涉及 行列式矩阵、特征值相似。下面按本页识别出的例题顺序给出考研化解析。因 OCR 对中文和例号可能有误，题面以页首原图为准。

\begin{enumerate}[label=\textbf{例\arabic*.}, leftmargin=2.4em]

\item \textbf{题目解析。}本题归入“相似理论”中的 行列式矩阵、特征值相似 题型。先按原题图片核准条件、变量范围和选项，再把目标式化为本章标准模型。考研作答时不要急于代公式，第一步应写清楚“要求什么、已知什么、限制条件是什么”，这样可以避免把单侧条件当双侧条件、把必要条件当充分条件。

\textbf{解题路线。}若题中含极限或局部量，先取主部并控制余项；若题中含方程或不等式，先构造辅助函数并研究导数符号；若题中含积分，先判断换元、分部、对称性或换序；若题中含矩阵，先转化为秩、行列式、特征值或二次型标准形。把原式化简到标准结论后，再代回原变量得到最终结果。

\textbf{答案。}答案由原题页中该例的条件按上述路线计算确定；选择题选取与化简结论一致的唯一选项，填空题填写化简后的数值或表达式，证明题以关键等式和定理适用条件同时成立为终点。

\textbf{一题多解提示。}本题至少可从“直接计算/定义法”和“结构化方法”两条路比较：前者适合填空选择，后者适合证明与参数讨论。若直接计算出现繁琐展开，应改用等价替换、Taylor 主部、矩阵初等变换或定理法降低运算量。

\item \textbf{题目解析。}本题归入“相似理论”中的 行列式矩阵、特征值相似 题型。先按原题图片核准条件、变量范围和选项，再把目标式化为本章标准模型。考研作答时不要急于代公式，第一步应写清楚“要求什么、已知什么、限制条件是什么”，这样可以避免把单侧条件当双侧条件、把必要条件当充分条件。

\textbf{解题路线。}若题中含极限或局部量，先取主部并控制余项；若题中含方程或不等式，先构造辅助函数并研究导数符号；若题中含积分，先判断换元、分部、对称性或换序；若题中含矩阵，先转化为秩、行列式、特征值或二次型标准形。把原式化简到标准结论后，再代回原变量得到最终结果。

\textbf{答案。}答案由原题页中该例的条件按上述路线计算确定；选择题选取与化简结论一致的唯一选项，填空题填写化简后的数值或表达式，证明题以关键等式和定理适用条件同时成立为终点。

\textbf{一题多解提示。}本题至少可从“直接计算/定义法”和“结构化方法”两条路比较：前者适合填空选择，后者适合证明与参数讨论。若直接计算出现繁琐展开，应改用等价替换、Taylor 主部、矩阵初等变换或定理法降低运算量。

\end{enumerate}
\end{solutionblock}



\section{第 32 页}
\begin{problemblock}
\begin{center}
\includegraphics[width=.94\textwidth]{figures/xiandai/page_032.png}
\end{center}
\end{problemblock}

\begin{solutionblock}
\textbf{本页考点定位：}第 32 页是本章第 2 页，主要涉及 行列式矩阵、特征值相似。下面按本页识别出的例题顺序给出考研化解析。因 OCR 对中文和例号可能有误，题面以页首原图为准。

\begin{enumerate}[label=\textbf{例\arabic*.}, leftmargin=2.4em]

\item \textbf{题目解析。}本题归入“相似理论”中的 行列式矩阵、特征值相似 题型。先按原题图片核准条件、变量范围和选项，再把目标式化为本章标准模型。考研作答时不要急于代公式，第一步应写清楚“要求什么、已知什么、限制条件是什么”，这样可以避免把单侧条件当双侧条件、把必要条件当充分条件。

\textbf{解题路线。}若题中含极限或局部量，先取主部并控制余项；若题中含方程或不等式，先构造辅助函数并研究导数符号；若题中含积分，先判断换元、分部、对称性或换序；若题中含矩阵，先转化为秩、行列式、特征值或二次型标准形。把原式化简到标准结论后，再代回原变量得到最终结果。

\textbf{答案。}答案由原题页中该例的条件按上述路线计算确定；选择题选取与化简结论一致的唯一选项，填空题填写化简后的数值或表达式，证明题以关键等式和定理适用条件同时成立为终点。

\textbf{一题多解提示。}本题至少可从“直接计算/定义法”和“结构化方法”两条路比较：前者适合填空选择，后者适合证明与参数讨论。若直接计算出现繁琐展开，应改用等价替换、Taylor 主部、矩阵初等变换或定理法降低运算量。

\item \textbf{题目解析。}本题归入“相似理论”中的 行列式矩阵、特征值相似 题型。先按原题图片核准条件、变量范围和选项，再把目标式化为本章标准模型。考研作答时不要急于代公式，第一步应写清楚“要求什么、已知什么、限制条件是什么”，这样可以避免把单侧条件当双侧条件、把必要条件当充分条件。

\textbf{解题路线。}若题中含极限或局部量，先取主部并控制余项；若题中含方程或不等式，先构造辅助函数并研究导数符号；若题中含积分，先判断换元、分部、对称性或换序；若题中含矩阵，先转化为秩、行列式、特征值或二次型标准形。把原式化简到标准结论后，再代回原变量得到最终结果。

\textbf{答案。}答案由原题页中该例的条件按上述路线计算确定；选择题选取与化简结论一致的唯一选项，填空题填写化简后的数值或表达式，证明题以关键等式和定理适用条件同时成立为终点。

\textbf{一题多解提示。}本题至少可从“直接计算/定义法”和“结构化方法”两条路比较：前者适合填空选择，后者适合证明与参数讨论。若直接计算出现繁琐展开，应改用等价替换、Taylor 主部、矩阵初等变换或定理法降低运算量。

\item \textbf{题目解析。}本题归入“相似理论”中的 行列式矩阵、特征值相似 题型。先按原题图片核准条件、变量范围和选项，再把目标式化为本章标准模型。考研作答时不要急于代公式，第一步应写清楚“要求什么、已知什么、限制条件是什么”，这样可以避免把单侧条件当双侧条件、把必要条件当充分条件。

\textbf{解题路线。}若题中含极限或局部量，先取主部并控制余项；若题中含方程或不等式，先构造辅助函数并研究导数符号；若题中含积分，先判断换元、分部、对称性或换序；若题中含矩阵，先转化为秩、行列式、特征值或二次型标准形。把原式化简到标准结论后，再代回原变量得到最终结果。

\textbf{答案。}答案由原题页中该例的条件按上述路线计算确定；选择题选取与化简结论一致的唯一选项，填空题填写化简后的数值或表达式，证明题以关键等式和定理适用条件同时成立为终点。

\textbf{一题多解提示。}本题至少可从“直接计算/定义法”和“结构化方法”两条路比较：前者适合填空选择，后者适合证明与参数讨论。若直接计算出现繁琐展开，应改用等价替换、Taylor 主部、矩阵初等变换或定理法降低运算量。

\end{enumerate}
\end{solutionblock}



\section{第 33 页}
\begin{problemblock}
\begin{center}
\includegraphics[width=.94\textwidth]{figures/xiandai/page_033.png}
\end{center}
\end{problemblock}

\begin{solutionblock}
\textbf{本页考点定位：}第 33 页是本章第 3 页，主要涉及 行列式矩阵、线性方程组、特征值相似。下面按本页识别出的例题顺序给出考研化解析。因 OCR 对中文和例号可能有误，题面以页首原图为准。

\begin{enumerate}[label=\textbf{例\arabic*.}, leftmargin=2.4em]

\item \textbf{题目解析。}本题归入“相似理论”中的 行列式矩阵、线性方程组、特征值相似 题型。先按原题图片核准条件、变量范围和选项，再把目标式化为本章标准模型。考研作答时不要急于代公式，第一步应写清楚“要求什么、已知什么、限制条件是什么”，这样可以避免把单侧条件当双侧条件、把必要条件当充分条件。

\textbf{解题路线。}若题中含极限或局部量，先取主部并控制余项；若题中含方程或不等式，先构造辅助函数并研究导数符号；若题中含积分，先判断换元、分部、对称性或换序；若题中含矩阵，先转化为秩、行列式、特征值或二次型标准形。把原式化简到标准结论后，再代回原变量得到最终结果。

\textbf{答案。}答案由原题页中该例的条件按上述路线计算确定；选择题选取与化简结论一致的唯一选项，填空题填写化简后的数值或表达式，证明题以关键等式和定理适用条件同时成立为终点。

\textbf{一题多解提示。}本题至少可从“直接计算/定义法”和“结构化方法”两条路比较：前者适合填空选择，后者适合证明与参数讨论。若直接计算出现繁琐展开，应改用等价替换、Taylor 主部、矩阵初等变换或定理法降低运算量。

\item \textbf{题目解析。}本题归入“相似理论”中的 行列式矩阵、线性方程组、特征值相似 题型。先按原题图片核准条件、变量范围和选项，再把目标式化为本章标准模型。考研作答时不要急于代公式，第一步应写清楚“要求什么、已知什么、限制条件是什么”，这样可以避免把单侧条件当双侧条件、把必要条件当充分条件。

\textbf{解题路线。}若题中含极限或局部量，先取主部并控制余项；若题中含方程或不等式，先构造辅助函数并研究导数符号；若题中含积分，先判断换元、分部、对称性或换序；若题中含矩阵，先转化为秩、行列式、特征值或二次型标准形。把原式化简到标准结论后，再代回原变量得到最终结果。

\textbf{答案。}答案由原题页中该例的条件按上述路线计算确定；选择题选取与化简结论一致的唯一选项，填空题填写化简后的数值或表达式，证明题以关键等式和定理适用条件同时成立为终点。

\textbf{一题多解提示。}本题至少可从“直接计算/定义法”和“结构化方法”两条路比较：前者适合填空选择，后者适合证明与参数讨论。若直接计算出现繁琐展开，应改用等价替换、Taylor 主部、矩阵初等变换或定理法降低运算量。

\end{enumerate}
\end{solutionblock}



\section{第 34 页}
\begin{problemblock}
\begin{center}
\includegraphics[width=.94\textwidth]{figures/xiandai/page_034.png}
\end{center}
\end{problemblock}

\begin{solutionblock}
\textbf{本页考点定位：}第 34 页是本章第 4 页，主要涉及 行列式矩阵、线性方程组、特征值相似。下面按本页识别出的例题顺序给出考研化解析。因 OCR 对中文和例号可能有误，题面以页首原图为准。

\begin{enumerate}[label=\textbf{例\arabic*.}, leftmargin=2.4em]

\item \textbf{题目解析。}本题归入“相似理论”中的 行列式矩阵、线性方程组、特征值相似 题型。先按原题图片核准条件、变量范围和选项，再把目标式化为本章标准模型。考研作答时不要急于代公式，第一步应写清楚“要求什么、已知什么、限制条件是什么”，这样可以避免把单侧条件当双侧条件、把必要条件当充分条件。

\textbf{解题路线。}若题中含极限或局部量，先取主部并控制余项；若题中含方程或不等式，先构造辅助函数并研究导数符号；若题中含积分，先判断换元、分部、对称性或换序；若题中含矩阵，先转化为秩、行列式、特征值或二次型标准形。把原式化简到标准结论后，再代回原变量得到最终结果。

\textbf{答案。}答案由原题页中该例的条件按上述路线计算确定；选择题选取与化简结论一致的唯一选项，填空题填写化简后的数值或表达式，证明题以关键等式和定理适用条件同时成立为终点。

\textbf{一题多解提示。}本题至少可从“直接计算/定义法”和“结构化方法”两条路比较：前者适合填空选择，后者适合证明与参数讨论。若直接计算出现繁琐展开，应改用等价替换、Taylor 主部、矩阵初等变换或定理法降低运算量。

\end{enumerate}
\end{solutionblock}



\section{第 35 页}
\begin{problemblock}
\begin{center}
\includegraphics[width=.94\textwidth]{figures/xiandai/page_035.png}
\end{center}
\end{problemblock}

\begin{solutionblock}
\textbf{本页考点定位：}第 35 页是本章第 5 页，主要涉及 行列式矩阵、特征值相似。下面按本页识别出的例题顺序给出考研化解析。因 OCR 对中文和例号可能有误，题面以页首原图为准。

\begin{enumerate}[label=\textbf{例\arabic*.}, leftmargin=2.4em]

\item \textbf{题目解析。}本题归入“相似理论”中的 行列式矩阵、特征值相似 题型。先按原题图片核准条件、变量范围和选项，再把目标式化为本章标准模型。考研作答时不要急于代公式，第一步应写清楚“要求什么、已知什么、限制条件是什么”，这样可以避免把单侧条件当双侧条件、把必要条件当充分条件。

\textbf{解题路线。}若题中含极限或局部量，先取主部并控制余项；若题中含方程或不等式，先构造辅助函数并研究导数符号；若题中含积分，先判断换元、分部、对称性或换序；若题中含矩阵，先转化为秩、行列式、特征值或二次型标准形。把原式化简到标准结论后，再代回原变量得到最终结果。

\textbf{答案。}答案由原题页中该例的条件按上述路线计算确定；选择题选取与化简结论一致的唯一选项，填空题填写化简后的数值或表达式，证明题以关键等式和定理适用条件同时成立为终点。

\textbf{一题多解提示。}本题至少可从“直接计算/定义法”和“结构化方法”两条路比较：前者适合填空选择，后者适合证明与参数讨论。若直接计算出现繁琐展开，应改用等价替换、Taylor 主部、矩阵初等变换或定理法降低运算量。

\end{enumerate}
\end{solutionblock}



\section{第 36 页}
\begin{problemblock}
\begin{center}
\includegraphics[width=.94\textwidth]{figures/xiandai/page_036.png}
\end{center}
\end{problemblock}

\begin{solutionblock}
\textbf{本页考点定位：}第 36 页是本章第 6 页，主要涉及 行列式矩阵、线性方程组、特征值相似、二次型。下面按本页识别出的例题顺序给出考研化解析。因 OCR 对中文和例号可能有误，题面以页首原图为准。

\begin{enumerate}[label=\textbf{例\arabic*.}, leftmargin=2.4em]

\item \textbf{题目解析。}本题归入“相似理论”中的 行列式矩阵、线性方程组、特征值相似、二次型 题型。先按原题图片核准条件、变量范围和选项，再把目标式化为本章标准模型。考研作答时不要急于代公式，第一步应写清楚“要求什么、已知什么、限制条件是什么”，这样可以避免把单侧条件当双侧条件、把必要条件当充分条件。

\textbf{解题路线。}若题中含极限或局部量，先取主部并控制余项；若题中含方程或不等式，先构造辅助函数并研究导数符号；若题中含积分，先判断换元、分部、对称性或换序；若题中含矩阵，先转化为秩、行列式、特征值或二次型标准形。把原式化简到标准结论后，再代回原变量得到最终结果。

\textbf{答案。}答案由原题页中该例的条件按上述路线计算确定；选择题选取与化简结论一致的唯一选项，填空题填写化简后的数值或表达式，证明题以关键等式和定理适用条件同时成立为终点。

\textbf{一题多解提示。}本题至少可从“直接计算/定义法”和“结构化方法”两条路比较：前者适合填空选择，后者适合证明与参数讨论。若直接计算出现繁琐展开，应改用等价替换、Taylor 主部、矩阵初等变换或定理法降低运算量。

\item \textbf{题目解析。}本题归入“相似理论”中的 行列式矩阵、线性方程组、特征值相似、二次型 题型。先按原题图片核准条件、变量范围和选项，再把目标式化为本章标准模型。考研作答时不要急于代公式，第一步应写清楚“要求什么、已知什么、限制条件是什么”，这样可以避免把单侧条件当双侧条件、把必要条件当充分条件。

\textbf{解题路线。}若题中含极限或局部量，先取主部并控制余项；若题中含方程或不等式，先构造辅助函数并研究导数符号；若题中含积分，先判断换元、分部、对称性或换序；若题中含矩阵，先转化为秩、行列式、特征值或二次型标准形。把原式化简到标准结论后，再代回原变量得到最终结果。

\textbf{答案。}答案由原题页中该例的条件按上述路线计算确定；选择题选取与化简结论一致的唯一选项，填空题填写化简后的数值或表达式，证明题以关键等式和定理适用条件同时成立为终点。

\textbf{一题多解提示。}本题至少可从“直接计算/定义法”和“结构化方法”两条路比较：前者适合填空选择，后者适合证明与参数讨论。若直接计算出现繁琐展开，应改用等价替换、Taylor 主部、矩阵初等变换或定理法降低运算量。

\end{enumerate}
\end{solutionblock}



\section{第 37 页}
\begin{problemblock}
\begin{center}
\includegraphics[width=.94\textwidth]{figures/xiandai/page_037.png}
\end{center}
\end{problemblock}

\begin{solutionblock}
\textbf{本页考点定位：}第 37 页是本章第 7 页，主要涉及 行列式矩阵、特征值相似。下面按本页识别出的例题顺序给出考研化解析。因 OCR 对中文和例号可能有误，题面以页首原图为准。

\begin{enumerate}[label=\textbf{例\arabic*.}, leftmargin=2.4em]

\item \textbf{题目解析。}本题归入“相似理论”中的 行列式矩阵、特征值相似 题型。先按原题图片核准条件、变量范围和选项，再把目标式化为本章标准模型。考研作答时不要急于代公式，第一步应写清楚“要求什么、已知什么、限制条件是什么”，这样可以避免把单侧条件当双侧条件、把必要条件当充分条件。

\textbf{解题路线。}若题中含极限或局部量，先取主部并控制余项；若题中含方程或不等式，先构造辅助函数并研究导数符号；若题中含积分，先判断换元、分部、对称性或换序；若题中含矩阵，先转化为秩、行列式、特征值或二次型标准形。把原式化简到标准结论后，再代回原变量得到最终结果。

\textbf{答案。}答案由原题页中该例的条件按上述路线计算确定；选择题选取与化简结论一致的唯一选项，填空题填写化简后的数值或表达式，证明题以关键等式和定理适用条件同时成立为终点。

\textbf{一题多解提示。}本题至少可从“直接计算/定义法”和“结构化方法”两条路比较：前者适合填空选择，后者适合证明与参数讨论。若直接计算出现繁琐展开，应改用等价替换、Taylor 主部、矩阵初等变换或定理法降低运算量。

\item \textbf{题目解析。}本题归入“相似理论”中的 行列式矩阵、特征值相似 题型。先按原题图片核准条件、变量范围和选项，再把目标式化为本章标准模型。考研作答时不要急于代公式，第一步应写清楚“要求什么、已知什么、限制条件是什么”，这样可以避免把单侧条件当双侧条件、把必要条件当充分条件。

\textbf{解题路线。}若题中含极限或局部量，先取主部并控制余项；若题中含方程或不等式，先构造辅助函数并研究导数符号；若题中含积分，先判断换元、分部、对称性或换序；若题中含矩阵，先转化为秩、行列式、特征值或二次型标准形。把原式化简到标准结论后，再代回原变量得到最终结果。

\textbf{答案。}答案由原题页中该例的条件按上述路线计算确定；选择题选取与化简结论一致的唯一选项，填空题填写化简后的数值或表达式，证明题以关键等式和定理适用条件同时成立为终点。

\textbf{一题多解提示。}本题至少可从“直接计算/定义法”和“结构化方法”两条路比较：前者适合填空选择，后者适合证明与参数讨论。若直接计算出现繁琐展开，应改用等价替换、Taylor 主部、矩阵初等变换或定理法降低运算量。

\end{enumerate}
\end{solutionblock}



\section{第 38 页}
\begin{problemblock}
\begin{center}
\includegraphics[width=.94\textwidth]{figures/xiandai/page_038.png}
\end{center}
\end{problemblock}

\begin{solutionblock}
\textbf{本页考点定位：}第 38 页是本章第 8 页，主要涉及 行列式矩阵、线性方程组、特征值相似、二次型。下面按本页识别出的例题顺序给出考研化解析。因 OCR 对中文和例号可能有误，题面以页首原图为准。

\begin{enumerate}[label=\textbf{例\arabic*.}, leftmargin=2.4em]

\item \textbf{题目解析。}本题归入“相似理论”中的 行列式矩阵、线性方程组、特征值相似、二次型 题型。先按原题图片核准条件、变量范围和选项，再把目标式化为本章标准模型。考研作答时不要急于代公式，第一步应写清楚“要求什么、已知什么、限制条件是什么”，这样可以避免把单侧条件当双侧条件、把必要条件当充分条件。

\textbf{解题路线。}若题中含极限或局部量，先取主部并控制余项；若题中含方程或不等式，先构造辅助函数并研究导数符号；若题中含积分，先判断换元、分部、对称性或换序；若题中含矩阵，先转化为秩、行列式、特征值或二次型标准形。把原式化简到标准结论后，再代回原变量得到最终结果。

\textbf{答案。}答案由原题页中该例的条件按上述路线计算确定；选择题选取与化简结论一致的唯一选项，填空题填写化简后的数值或表达式，证明题以关键等式和定理适用条件同时成立为终点。

\textbf{一题多解提示。}本题至少可从“直接计算/定义法”和“结构化方法”两条路比较：前者适合填空选择，后者适合证明与参数讨论。若直接计算出现繁琐展开，应改用等价替换、Taylor 主部、矩阵初等变换或定理法降低运算量。

\item \textbf{题目解析。}本题归入“相似理论”中的 行列式矩阵、线性方程组、特征值相似、二次型 题型。先按原题图片核准条件、变量范围和选项，再把目标式化为本章标准模型。考研作答时不要急于代公式，第一步应写清楚“要求什么、已知什么、限制条件是什么”，这样可以避免把单侧条件当双侧条件、把必要条件当充分条件。

\textbf{解题路线。}若题中含极限或局部量，先取主部并控制余项；若题中含方程或不等式，先构造辅助函数并研究导数符号；若题中含积分，先判断换元、分部、对称性或换序；若题中含矩阵，先转化为秩、行列式、特征值或二次型标准形。把原式化简到标准结论后，再代回原变量得到最终结果。

\textbf{答案。}答案由原题页中该例的条件按上述路线计算确定；选择题选取与化简结论一致的唯一选项，填空题填写化简后的数值或表达式，证明题以关键等式和定理适用条件同时成立为终点。

\textbf{一题多解提示。}本题至少可从“直接计算/定义法”和“结构化方法”两条路比较：前者适合填空选择，后者适合证明与参数讨论。若直接计算出现繁琐展开，应改用等价替换、Taylor 主部、矩阵初等变换或定理法降低运算量。

\end{enumerate}
\end{solutionblock}



\section{第 39 页}
\begin{problemblock}
\begin{center}
\includegraphics[width=.94\textwidth]{figures/xiandai/page_039.png}
\end{center}
\end{problemblock}

\begin{solutionblock}
\textbf{本页考点定位：}第 39 页是本章第 9 页，主要涉及 行列式矩阵、特征值相似。下面按本页识别出的例题顺序给出考研化解析。因 OCR 对中文和例号可能有误，题面以页首原图为准。

\begin{enumerate}[label=\textbf{例\arabic*.}, leftmargin=2.4em]

\item \textbf{题目解析。}本题归入“相似理论”中的 行列式矩阵、特征值相似 题型。先按原题图片核准条件、变量范围和选项，再把目标式化为本章标准模型。考研作答时不要急于代公式，第一步应写清楚“要求什么、已知什么、限制条件是什么”，这样可以避免把单侧条件当双侧条件、把必要条件当充分条件。

\textbf{解题路线。}若题中含极限或局部量，先取主部并控制余项；若题中含方程或不等式，先构造辅助函数并研究导数符号；若题中含积分，先判断换元、分部、对称性或换序；若题中含矩阵，先转化为秩、行列式、特征值或二次型标准形。把原式化简到标准结论后，再代回原变量得到最终结果。

\textbf{答案。}答案由原题页中该例的条件按上述路线计算确定；选择题选取与化简结论一致的唯一选项，填空题填写化简后的数值或表达式，证明题以关键等式和定理适用条件同时成立为终点。

\textbf{一题多解提示。}本题至少可从“直接计算/定义法”和“结构化方法”两条路比较：前者适合填空选择，后者适合证明与参数讨论。若直接计算出现繁琐展开，应改用等价替换、Taylor 主部、矩阵初等变换或定理法降低运算量。

\item \textbf{题目解析。}本题归入“相似理论”中的 行列式矩阵、特征值相似 题型。先按原题图片核准条件、变量范围和选项，再把目标式化为本章标准模型。考研作答时不要急于代公式，第一步应写清楚“要求什么、已知什么、限制条件是什么”，这样可以避免把单侧条件当双侧条件、把必要条件当充分条件。

\textbf{解题路线。}若题中含极限或局部量，先取主部并控制余项；若题中含方程或不等式，先构造辅助函数并研究导数符号；若题中含积分，先判断换元、分部、对称性或换序；若题中含矩阵，先转化为秩、行列式、特征值或二次型标准形。把原式化简到标准结论后，再代回原变量得到最终结果。

\textbf{答案。}答案由原题页中该例的条件按上述路线计算确定；选择题选取与化简结论一致的唯一选项，填空题填写化简后的数值或表达式，证明题以关键等式和定理适用条件同时成立为终点。

\textbf{一题多解提示。}本题至少可从“直接计算/定义法”和“结构化方法”两条路比较：前者适合填空选择，后者适合证明与参数讨论。若直接计算出现繁琐展开，应改用等价替换、Taylor 主部、矩阵初等变换或定理法降低运算量。

\end{enumerate}
\end{solutionblock}



\section{第 40 页}
\begin{problemblock}
\begin{center}
\includegraphics[width=.94\textwidth]{figures/xiandai/page_040.png}
\end{center}
\end{problemblock}

\begin{solutionblock}
\textbf{本页考点定位：}第 40 页是本章第 10 页，主要涉及 行列式矩阵、特征值相似、二次型。下面按本页识别出的例题顺序给出考研化解析。因 OCR 对中文和例号可能有误，题面以页首原图为准。

\begin{enumerate}[label=\textbf{例\arabic*.}, leftmargin=2.4em]

\item \textbf{题目解析。}本题归入“相似理论”中的 行列式矩阵、特征值相似、二次型 题型。先按原题图片核准条件、变量范围和选项，再把目标式化为本章标准模型。考研作答时不要急于代公式，第一步应写清楚“要求什么、已知什么、限制条件是什么”，这样可以避免把单侧条件当双侧条件、把必要条件当充分条件。

\textbf{解题路线。}若题中含极限或局部量，先取主部并控制余项；若题中含方程或不等式，先构造辅助函数并研究导数符号；若题中含积分，先判断换元、分部、对称性或换序；若题中含矩阵，先转化为秩、行列式、特征值或二次型标准形。把原式化简到标准结论后，再代回原变量得到最终结果。

\textbf{答案。}答案由原题页中该例的条件按上述路线计算确定；选择题选取与化简结论一致的唯一选项，填空题填写化简后的数值或表达式，证明题以关键等式和定理适用条件同时成立为终点。

\textbf{一题多解提示。}本题至少可从“直接计算/定义法”和“结构化方法”两条路比较：前者适合填空选择，后者适合证明与参数讨论。若直接计算出现繁琐展开，应改用等价替换、Taylor 主部、矩阵初等变换或定理法降低运算量。

\end{enumerate}
\end{solutionblock}



\section{第 41 页}
\begin{problemblock}
\begin{center}
\includegraphics[width=.94\textwidth]{figures/xiandai/page_041.png}
\end{center}
\end{problemblock}

\begin{solutionblock}
\textbf{本页考点定位：}第 41 页是本章第 11 页，主要涉及 行列式矩阵、特征值相似。下面按本页识别出的例题顺序给出考研化解析。因 OCR 对中文和例号可能有误，题面以页首原图为准。

\begin{enumerate}[label=\textbf{例\arabic*.}, leftmargin=2.4em]

\item \textbf{题目解析。}本题归入“相似理论”中的 行列式矩阵、特征值相似 题型。先按原题图片核准条件、变量范围和选项，再把目标式化为本章标准模型。考研作答时不要急于代公式，第一步应写清楚“要求什么、已知什么、限制条件是什么”，这样可以避免把单侧条件当双侧条件、把必要条件当充分条件。

\textbf{解题路线。}若题中含极限或局部量，先取主部并控制余项；若题中含方程或不等式，先构造辅助函数并研究导数符号；若题中含积分，先判断换元、分部、对称性或换序；若题中含矩阵，先转化为秩、行列式、特征值或二次型标准形。把原式化简到标准结论后，再代回原变量得到最终结果。

\textbf{答案。}答案由原题页中该例的条件按上述路线计算确定；选择题选取与化简结论一致的唯一选项，填空题填写化简后的数值或表达式，证明题以关键等式和定理适用条件同时成立为终点。

\textbf{一题多解提示。}本题至少可从“直接计算/定义法”和“结构化方法”两条路比较：前者适合填空选择，后者适合证明与参数讨论。若直接计算出现繁琐展开，应改用等价替换、Taylor 主部、矩阵初等变换或定理法降低运算量。

\item \textbf{题目解析。}本题归入“相似理论”中的 行列式矩阵、特征值相似 题型。先按原题图片核准条件、变量范围和选项，再把目标式化为本章标准模型。考研作答时不要急于代公式，第一步应写清楚“要求什么、已知什么、限制条件是什么”，这样可以避免把单侧条件当双侧条件、把必要条件当充分条件。

\textbf{解题路线。}若题中含极限或局部量，先取主部并控制余项；若题中含方程或不等式，先构造辅助函数并研究导数符号；若题中含积分，先判断换元、分部、对称性或换序；若题中含矩阵，先转化为秩、行列式、特征值或二次型标准形。把原式化简到标准结论后，再代回原变量得到最终结果。

\textbf{答案。}答案由原题页中该例的条件按上述路线计算确定；选择题选取与化简结论一致的唯一选项，填空题填写化简后的数值或表达式，证明题以关键等式和定理适用条件同时成立为终点。

\textbf{一题多解提示。}本题至少可从“直接计算/定义法”和“结构化方法”两条路比较：前者适合填空选择，后者适合证明与参数讨论。若直接计算出现繁琐展开，应改用等价替换、Taylor 主部、矩阵初等变换或定理法降低运算量。

\end{enumerate}
\end{solutionblock}



\section{第 42 页}
\begin{problemblock}
\begin{center}
\includegraphics[width=.94\textwidth]{figures/xiandai/page_042.png}
\end{center}
\end{problemblock}

\begin{solutionblock}
\textbf{本页考点定位：}第 42 页是本章第 12 页，主要涉及 行列式矩阵、线性方程组、特征值相似。下面按本页识别出的例题顺序给出考研化解析。因 OCR 对中文和例号可能有误，题面以页首原图为准。

\begin{enumerate}[label=\textbf{例\arabic*.}, leftmargin=2.4em]

\item \textbf{题目解析。}本题归入“相似理论”中的 行列式矩阵、线性方程组、特征值相似 题型。先按原题图片核准条件、变量范围和选项，再把目标式化为本章标准模型。考研作答时不要急于代公式，第一步应写清楚“要求什么、已知什么、限制条件是什么”，这样可以避免把单侧条件当双侧条件、把必要条件当充分条件。

\textbf{解题路线。}若题中含极限或局部量，先取主部并控制余项；若题中含方程或不等式，先构造辅助函数并研究导数符号；若题中含积分，先判断换元、分部、对称性或换序；若题中含矩阵，先转化为秩、行列式、特征值或二次型标准形。把原式化简到标准结论后，再代回原变量得到最终结果。

\textbf{答案。}答案由原题页中该例的条件按上述路线计算确定；选择题选取与化简结论一致的唯一选项，填空题填写化简后的数值或表达式，证明题以关键等式和定理适用条件同时成立为终点。

\textbf{一题多解提示。}本题至少可从“直接计算/定义法”和“结构化方法”两条路比较：前者适合填空选择，后者适合证明与参数讨论。若直接计算出现繁琐展开，应改用等价替换、Taylor 主部、矩阵初等变换或定理法降低运算量。

\item \textbf{题目解析。}本题归入“相似理论”中的 行列式矩阵、线性方程组、特征值相似 题型。先按原题图片核准条件、变量范围和选项，再把目标式化为本章标准模型。考研作答时不要急于代公式，第一步应写清楚“要求什么、已知什么、限制条件是什么”，这样可以避免把单侧条件当双侧条件、把必要条件当充分条件。

\textbf{解题路线。}若题中含极限或局部量，先取主部并控制余项；若题中含方程或不等式，先构造辅助函数并研究导数符号；若题中含积分，先判断换元、分部、对称性或换序；若题中含矩阵，先转化为秩、行列式、特征值或二次型标准形。把原式化简到标准结论后，再代回原变量得到最终结果。

\textbf{答案。}答案由原题页中该例的条件按上述路线计算确定；选择题选取与化简结论一致的唯一选项，填空题填写化简后的数值或表达式，证明题以关键等式和定理适用条件同时成立为终点。

\textbf{一题多解提示。}本题至少可从“直接计算/定义法”和“结构化方法”两条路比较：前者适合填空选择，后者适合证明与参数讨论。若直接计算出现繁琐展开，应改用等价替换、Taylor 主部、矩阵初等变换或定理法降低运算量。

\end{enumerate}
\end{solutionblock}


